\begin{claim}{B.27}{Proved}{Complete executed punctured-torus catalogues through four punctures}
The extended mapping class groups of a torus with $\nu=1,2,3,4$
unlabelled punctures have, respectively,
\[
                         13,\quad40,\quad38,\quad115
\]
conjugacy classes of finite subgroups.  The files
\texttt{type\_1\_\nu\_affine\_catalogue.csv}, together with their
\texttt{affine\_sections.json} companions and
\texttt{torus\_affine\_normal\_forms.json}, list every class once.
They specify each group by its complete finite set of rational affine
transformations, its punctures, every full and nontrivial partial section
class, all extension incidences, $\Phi$ and its group orbits and stabilizers,
the centralizers, and $\mu$ on every section and puncture-stabilizer element.
Thus the torus portion of B(f) is complete for every hyperbolic type in the
requested range.
\end{claim}
\begin{proof}
We give the finite classification underlying the program
\texttt{catalogue\_tori.py}; this also specifies the tables without relying
on abstract finite-group names.  By B.14 extend over the punctures and put
the action in affine form on $V/\Lambda=\R^2/\Z^2$.  Write
\[
 1\longrightarrow T\longrightarrow G\longrightarrow H\longrightarrow1,
 \qquad T\subset V/\Lambda,\quad H\subset\mathrm{GL}_2(\Z).
                                                        \tag{B.27.1}
\]
Here $T$ consists exactly of the translations in $G$.  It acts freely on
$V/\Lambda$, so each $T$-orbit on the puncture set has $|T|$ elements.
Consequently $|T|\mid\nu$ and in particular $|T|\le4$.  The complete list
of subgroups of $V/\Lambda$ of order at most four consists of the trivial
subgroup, the three cyclic subgroups of order two, the four cyclic
subgroups of order three, the six cyclic subgroups of order four, and
$(\frac12\Lambda)/\Lambda$.  Indeed a group of order at most four is
cyclic or $C_2^2$; a cyclic subgroup is generated by a point of that exact
order, and division by the number of generators gives precisely these
counts.  All fifteen groups lie in $(\frac1{12}\Lambda)/\Lambda$.
The thirteen possibilities for $H$, up to integral conjugacy, are exactly
those proved in B.25.  For each we retain only the $H$-invariant $T$.

We now enumerate the affine extensions, including the actions with no
common fixed point.  If $H$ contains a nonidentity positive element,
choose a generator $r$ of its cyclic positive subgroup.  Every affine
map with linear part $r$ has a fixed point, since $I-r$ is invertible
on $V$ and is therefore onto on $V/\Lambda$.  Translation of that point
to the origin makes this generator $(r,0)$.  A positive cyclic $H$
then gives exactly $G=\langle T,(r,0)\rangle$.  For a dihedral $H$
choose a negative generator $f$ and write its lift as $(f,b)$.  The two
relations in $H$ give exactly
\[
       (I-r)b\in T,\qquad (I+f)b\in T.                 \tag{B.27.2}
\]
Conversely these equations make the subgroup generated by
$T,(r,0),(f,b)$ a finite group with kernel $T$ and quotient $H$:
every word reduces to a translation followed by $r^j$ or $r^jf$, and
its linear part distinguishes these cosets.  Values of $b$ modulo $T$
are enough.  The finite set $(I-r)^{-1}T$ therefore supplies every
candidate, with no bound on words or finite quotients involved.
Translations of the origin preserving this normal form are exactly
$a$ with $(I-r)a\in T$; they change $b$ by $(I-f)a$ modulo $T$.

The remaining nontrivial linear groups have $H=\langle f\rangle$
of order two.  Put $W=(V/\Lambda)/T$.  Their extensions are represented
by $b\in W$ with $(I+f)b=0$, modulo $(I-f)W$.  Every such class has a
representative in $W[2]$.  To verify this assertion directly, choose a
lift $v\in V$ of $b$ for the enlarged lattice
$\Lambda_T=\{v\in V:v\bmod\Lambda\in T\}$.
Then $(I+f)v\in\Lambda_T$.  Write
$v=v_++v_-$ in the real $+1$ and $-1$ eigenspaces of $f$.
Subtracting $(I-f)(v_-/2)$ removes $v_-$, and the remaining $v_+$
satisfies $2v_+\in\Lambda_T$.  Thus the program uses precisely
\[
              2b\in T,\qquad (I+f)b\in T.             \tag{B.27.3}
\]
For two candidates $b,b'$ the translation-conjugacy test is also exact.
Let $w$ be a primitive integral row vector satisfying $wf=w$.
The image of $I-f$ on $V/\Lambda$ is the connected circle $\ker w$:
its real image is that eigenspace, and multiplication by any nonzero
integer is onto on a circle.  Hence the two candidates are equivalent
precisely when $w(b-b')\in w(T)$ in $\R/\Z$.
This includes both reflections and fixed-point-free glide reflections.

Integral changes of basis are handled without truncating an infinite
matrix group.  For $H=1$ or $\{\pm I\}$, use the generators $S,R,F$
of B.25 on the complete fifteen-element list of possible $T$'s.
Their orbit calculation is the full $\mathrm{GL}_2(\Z)$-conjugacy
calculation.  For a fixed $T$, a spanning tree of this finite orbit
and Schreier's formula give generators for its entire integral
stabilizer.  Explicitly, with a chosen representative $B_U$ carrying
$T$ to $U$, the generator associated to a matrix $A\in\{S,R,F,S^{-1},R^{-1}\}$
and the edge $U\mapsto AU$ is
\[
                       B_{AU}^{-1}AB_U.                \tag{B.27.4}
\]
For every other $H$ its integral normalizer is finite.  Here is an
exhaustiveness check for the small matrix enumeration in the code.
For a reflection conjugate to $\operatorname{diag}(-1,1)$, a commuting
integral unimodular matrix is diagonal with entries $\pm1$.
For a reflection conjugate to
$\left(\begin{smallmatrix}0&-1\\-1&0\end{smallmatrix}\right)$,
the equation is $A=\left(\begin{smallmatrix}a&b\\b&a\end{smallmatrix}\right)$,
with $(a-b)(a+b)=\pm1$, giving four solutions.  A normalizer of the
four-group generated by this reflection and $-I$ may also exchange
the two eigenspaces; the same calculation gives eight solutions.
A group with a rotation of order four has its normalizer in the eight
signed permutation matrices, since a commuting matrix has determinant
$a^2+b^2=1$.  A group with rotation of order three or six has its
normalizer in the twelve matrices for the hexagonal lattice: the
commuting equation gives the norm $a^2+ab+b^2=1$, and one reflection
supplies the other coset.  Testing these finite sets against $H$
gives the full normalizer.  Every resulting matrix has entries of
absolute value at most one, so the code's box of entries from $-2$
to $2$ contains every solution and introduces none outside the tested
normalizer equations.  Applying these matrices and the origin changes
above fuses all and only affine-conjugate candidates.
The result is $93$ affine action classes with $|T|\le4$ before choosing
punctures; some admit no invariant set of a given required size.

We explain why decorating these affine actions gives topological
conjugacy, including all permutations of the singular strata.
B.14 proves that a topological conjugacy between finite affine actions
can be replaced by an affine conjugacy.  The replacement also induces
the same permutation on the strata used here.  In detail write a lift
of the conjugacy as $h(x)=Ax+\delta(x)$ with $\delta$ periodic, and
put $\bar\delta=\int_{V/\Lambda}\delta$.  The equivariance equations
and invariance of Haar measure show that $x\mapsto Ax+\bar\delta$
has the same equivariance.  The straight interpolation of these two
lifts is an equivariant homotopy.  An isolated fixed point of a
nontrivial positive subgroup stays constant under this homotopy.
Each reflection circle stays in the same fixed-set component.
The degree on every such circle is fixed by the homology map $A$;
these circles are essential, since they lift to parallel lines.
Together with the images of the isolated vertices this determines
the permutation of the intervening open mirror intervals.  Thus the
affine normalizer gives exactly the possible stratum permutations.

To compute them, let $\mathcal V$ be the union of the fixed points of
all positive nontranslation elements.  It is a finite set.  The fixed
circles of the reflections, cut at $\mathcal V$, are open mirror
intervals or uncut mirror circles.  Take the $G$-orbits of the vertices
and of these mirror components; call their numbers $v$ and $e$.
There is one regular interior stratum in the quotient, since its
interior after removal of finitely many cone points is connected.
A vertex orbit may be removed once, a mirror-component orbit may
contain any nonnegative number of puncture orbits, and the regular
interior may contain any nonnegative number.  Thus the decoration is
\[
 (d_1,\ldots,d_v;c_1,\ldots,c_e;a),\qquad
 d_i\in\{0,1\},\ c_j,a\ge0,
\]
subject to the exact equation
\[
      \sum_i d_i|V_i|+\sum_j c_j\frac{|G|}{2}+a|G|=\nu.
                                                        \tag{B.27.5}
\]
Divide these finitely many decorations by the stratum permutations
of the affine normalizer computed above.  Counts are sufficient on
an open mirror interval or circle: an increasing interval map, or a
circle map of the prescribed degree, moves any two unlabelled finite
sets of the same size into each other.  The moves extend through
small strips in the quotient.  In the regular interior use disjoint
moving discs; its configuration space of finitely many unlabelled
points is connected.  Such quotient isotopies lift through the finite
orbifold covering.  They preserve all cone and corner charts and
therefore give equivariant homeomorphisms upstairs.  Conversely a
conjugacy preserves the displayed orbit sizes and the incidence of
these strata.  This proves completeness and injectivity of the
decorated classification.

All computations use exact integers on
$Q=(\frac1{144}\Z^2)/\Z^2$.  We record why this finite mesh loses no
component or conjugacy.  Coordinates of $T$ have denominator dividing
its order $k\le4$.  For the rotations of orders $2,3,4,6$, the inverses
of $I-r$ have denominators $2,3,2,1$, respectively.  Thus all solutions
needed for (B.27.2), all normalized origin changes, and all isolated
fixed points have denominators dividing $2k$ or $3k$.
The representatives (B.27.3) have denominator dividing $2k$.
For a reflection normalizer, the origin equation is a rank-one
integral equation with nonzero Smith invariant $1$ or $2$; every
connected solution component has a representative with denominator
dividing $4k$.  The integer $144$ is divisible by each of these
integers for $1\le k\le4$.  The code separately checks the inverse
matrix and particular-solution denominators for every positive affine
element.  On each reflection circle, moving by its primitive integral
direction divided by $144$ connects consecutive sampled points.
Every isolated intersection and every selected puncture is retained
as a vertex.  This graph is a subdivision of the actual fixed-circle
arrangement, so its components after removing punctures are exactly
the ovals and lines of the action.  The mesh has enough unused
representatives on each eligible stratum to place the at most four
punctures; the executed program explicitly verifies the prescribed
orbit counts and total cardinality for every decoration.

For each unremoved vertex $p$, enumerate all subgroups of its finite
cyclic or dihedral stabilizer that contain a positive nonidentity
element.  Each gives the section class with fixed point $p$.
For each reflection enumerate the components of its fixed circles
after deleting the punctures.  B.3--B.4 identify exactly these data
with all nontrivial partial-section classes, and identify their
centralizers as $1$ at points and lines and $\Z$ at ovals.
Containment of image subgroups and reverse containment of their fixed
components give all extension relations.  Applying every affine
matrix computes the group action on this poset and hence $\Phi$,
its orbit decomposition and every stabilizer.  At a positive fixed
point the local matrix is its actual linear part.  Its trace
$-2,-1,0,1$ gives order $2,3,4,6$, and the sign of
$\det(e_1,Ae_1)$ gives the sign of the angle except for a half-turn.
Reflections have multiplier $1/2$.  The program checks the order
and power identities on every element of every section.  It retains
rather than deletes the reflection section classes with several
maximal extensions.

Finally a nonidentity finite affine map acting trivially in the
punctured-torus mapping class group would have trivial homology
matrix and hence be a nonidentity translation.  It would also fix
each puncture, contradicting freeness of translations.  Thus each
action gives a faithful finite mapping class subgroup.  The realization
and uniqueness statements C.5--C.6 identify their conjugacy classes
with the topological action classes just enumerated.  The computation
returns the four counts in the statement.  For $\nu=1$, the multiset
of group orders, positive orders, full splitting counts and partial
splitting counts agrees, class for class, with the independent
amalgam calculation of B.25.  A first complete run with denominator
$72$ and the final run with denominator $144$ also give the same
four counts; the proof of completeness uses the latter mesh.
\end{proof}
